% orbit_ellipse: the vocabulary of an elliptical orbit (Ch 17, fig:orbitellipse)
% Semi-major axis a = 2.3, eccentricity e = 0.6, so b = a sqrt(1-e^2) = 1.84 and
% the foci sit at x = +-ae = +-1.38.  The Sun is at the right focus, so the
% perihelion (x = a) is at distance a(1-e) and the aphelion at a(1+e).
% The point Q at parameter angle 120 deg shows r1 + r2 = 2a.
\begin{tikzpicture}
  \pgfmathsetmacro\a{2.3}
  \pgfmathsetmacro\e{0.6}
  \pgfmathsetmacro\b{\a*sqrt(1-\e*\e)}
  \pgfmathsetmacro\c{\a*\e}
  \coordinate (C) at (0,0);
  \coordinate (S) at (\c,0);
  \coordinate (F) at (-\c,0);
  \coordinate (Q) at ({\a*cos(115)},{\b*sin(115)});
  % major axis
  \draw[guide] (-\a-0.3,0) -- (\a+0.3,0);
  % the two focal distances to a point Q on the orbit
  \draw[cTeal,thin line] (F) -- (Q) node[pos=0.5,anchor=east,inner sep=2pt] {$r_2$};
  \draw[cTeal,thin line] (S) -- (Q) node[pos=0.6,anchor=south west,inner sep=1pt] {$r_1$};
  \node[cTeal,anchor=south east,inner sep=2pt] at (Q) {$r_1+r_2=2a$};
  % the orbit
  \draw[cPrim,curve] (C) ellipse[x radius=\a,y radius=\b];
  \node[point mass] at (Q) {};
  % semi-minor axis
  \draw[cSteel,thin line] (C) -- (0,-\b) node[pos=0.5,anchor=west,inner sep=1.5pt] {$b$};
  % centre to focus
  \draw[dim] (0,0.32) -- (\c,0.32) node[midway,anchor=north,inner sep=1.5pt] {$ae$};
  \draw[dimtick] (\c,0) -- (\c,0.45);
  % semi-major axis, above the orbit
  \draw[dim] (0,\b+0.35) -- (\a,\b+0.35) node[midway,anchor=south,inner sep=1pt] {$a$};
  \draw[dimtick] (0,\b) -- (0,\b+0.45);
  \draw[dimtick] (\a,0) -- (\a,\b+0.45);
  % perihelion and aphelion distances, below the orbit
  \draw[dimtick] (\c,0) -- (\c,-\b-0.45);
  \draw[dimtick] (\a,0) -- (\a,-\b-0.45);
  \draw[dimtick] (-\a,0) -- (-\a,-\b-0.45);
  \draw[dim] (\c,-\b-0.35) -- (\a,-\b-0.35);
  \draw[dim] (-\a,-\b-0.35) -- (\c,-\b-0.35);
  \node[cSteel,anchor=north,inner sep=1pt] at ({(\c+\a)/2},-\b-0.5) {$r_p=a(1-e)$};
  \node[cSteel,anchor=north,inner sep=1pt] at ({(\c-\a)/2},-\b-0.5) {$r_a=a(1+e)$};
  % foci and centre
  \node[point mass] at (C) {};
  \node[anchor=north,inner sep=2pt] at (C) {centre};
  \filldraw[fill=cSun,draw=cAcc,thin line] (S) circle[radius=0.1];
  \node[anchor=north,inner sep=3pt] at (S) {Sun};
  \node[point mass,fill=cGrey] at (F) {};
  \node[anchor=north,inner sep=3pt,align=center,text=cSteel] at (F) {empty\\ focus};
  % apsides
  \node[anchor=west,inner sep=4pt] at (\a,0.25) {perihelion};
  \node[anchor=east,inner sep=3pt] at (-\a,0.25) {aphelion};
\end{tikzpicture}
